\section{Datos y método}

\subsection{Unidad de análisis y periodo}

La investigación adopta un diseño cuantitativo, no experimental y longitudinal
comparativo. La unidad de análisis corresponde a los departamentos del Perú y
la muestra comprende 24 unidades territoriales. La Provincia Constitucional
del Callao fue excluida debido a que la serie departamental utilizada para
construir el producto bruto interno real no presenta un archivo separado para
dicha jurisdicción.

El periodo de análisis comprende los años 2007 y 2019. Esta ventana responde a
la necesidad de trabajar con información comparable y metodológicamente
homogénea para el Índice de Desarrollo Humano, el producto bruto interno real y
la población departamental. Asimismo, 2007 coincide con el año base de la serie
de producción regional a precios constantes, mientras que 2019 constituye el
último año disponible en el panel principal de desarrollo humano empleado en
la investigación.

El trabajo constituye una adaptación metodológica del estudio de
\textcite{duran2024}. La réplica conserva el análisis distributivo, la
econometría espacial y la estimación de parámetros locales, pero ajusta el
periodo, las unidades territoriales y las variables a la disponibilidad de
información regional del Perú.

\subsection{Fuentes de información}

El producto bruto interno departamental se obtuvo de la serie de producto
bruto interno por departamentos según actividades económicas del Instituto
Nacional de Estadística e Informática
\parencite{inei_pbi}. Se utilizó el Valor Agregado Bruto a precios constantes
de 2007, expresado en miles de soles. El valor departamental fue construido a
partir de la suma de las actividades económicas reportadas en los archivos
correspondientes a cada territorio.

La población departamental procede de las estimaciones y proyecciones de
población del INEI \parencite{inei_poblacion}. Se empleó la población total
estimada al 30 de junio de cada año. Esta información permitió transformar el
Valor Agregado Bruto real en un indicador por habitante y evitar que el tamaño
poblacional determinara directamente la comparación entre departamentos.

El Índice de Desarrollo Humano fue obtenido del panel departamental
2003--2019 publicado por el Programa de las Naciones Unidas para el Desarrollo
\parencite{pnud_idh,pnud2019}. Se seleccionaron las observaciones
departamentales correspondientes a 2007 y 2019 mediante los códigos territoriales
del archivo original.

La información geográfica proviene del archivo de límites departamentales
publicado en la Plataforma Nacional de Datos Abiertos
\parencite{limites_departamentales}. El archivo contiene los polígonos
departamentales necesarios para elaborar los mapas y construir las matrices de
pesos espaciales.

\begin{table}[H]
\centering
\caption{Fuentes y variables utilizadas}
\label{tab:fuentes_variables}
\small
\begin{threeparttable}
\begin{tabular}{p{3.1cm}p{4.1cm}p{5.9cm}}
\toprule
\textbf{Componente} &
\textbf{Fuente} &
\textbf{Uso en el estudio} \\
\midrule
Valor Agregado Bruto real &
Instituto Nacional de Estadística e Informática &
Construcción del producto bruto interno real departamental a precios
constantes de 2007. \\

Población &
Instituto Nacional de Estadística e Informática &
Cálculo del producto bruto interno real por habitante para cada departamento. \\

Índice de Desarrollo Humano &
Programa de las Naciones Unidas para el Desarrollo &
Medición multidimensional del desarrollo departamental en 2007 y 2019. \\

Límites departamentales &
IGN y Plataforma Nacional de Datos Abiertos &
Elaboración de mapas y matrices de proximidad geográfica. \\
\bottomrule
\end{tabular}
\begin{tablenotes}
\footnotesize
\item \textit{Nota.} El análisis comprende 24 departamentos. Callao fue
excluido por la falta de una serie separada y compatible de producción
departamental.
\end{tablenotes}
\end{threeparttable}
\end{table}

\subsection{Construcción de variables}

El producto bruto interno real per cápita del departamento $i$ en el año $t$
se calculó mediante:

\begin{equation}
PBIpc_{it}
=
\frac{VAB_{it}}{Población_{it}},
\label{eq:pbipc}
\end{equation}

donde $VAB_{it}$ representa el Valor Agregado Bruto a precios constantes de
2007 y $Población_{it}$ corresponde a la población departamental estimada.

Para facilitar la comparación entre territorios y controlar la evolución
agregada de la economía peruana, se construyó un indicador relativo:

\begin{equation}
PBIpc^{rel}_{it}
=
\frac{PBIpc_{it}}{PBIpc^{PER}_{t}},
\label{eq:pbipc_relativo}
\end{equation}

donde $PBIpc^{PER}_{t}$ es el producto bruto interno real per cápita nacional.
Posteriormente se aplicó el logaritmo natural:

\begin{equation}
x_{it}
=
\ln\left(PBIpc^{rel}_{it}\right).
\label{eq:log_pbipc}
\end{equation}

El cambio acumulado del indicador económico se definió como:

\begin{equation}
\Delta x_i
=
x_{i,2019}-x_{i,2007}.
\label{eq:cambio_pbi}
\end{equation}

Para el desarrollo humano se utilizó directamente el valor departamental del
IDH, cuya escala se encuentra comprendida entre cero y uno. Esta decisión
constituye una adaptación respecto del estudio original y evita realizar una
normalización adicional sobre un indicador que ya posee una escala común. El
cambio acumulado se calculó mediante:

\begin{equation}
\Delta IDH_i
=
IDH_{i,2019}-IDH_{i,2007}.
\label{eq:cambio_idh}
\end{equation}

\begin{table}[H]
\centering
\caption{Definición de las variables principales}
\label{tab:definicion_variables}
\small
\begin{threeparttable}
\begin{tabular}{p{3.5cm}p{3.2cm}p{6.2cm}}
\toprule
\textbf{Variable} &
\textbf{Notación} &
\textbf{Definición} \\
\midrule
IDH inicial &
$IDH_{i,2007}$ &
Índice de Desarrollo Humano departamental en 2007. \\

Cambio del IDH &
$\Delta IDH_i$ &
Diferencia entre el IDH de 2019 y el IDH de 2007. \\

PBI per cápita relativo inicial &
$x_{i,2007}$ &
Logaritmo del cociente entre el PBI real per cápita departamental y el
nacional en 2007. \\

Cambio del PBI per cápita relativo &
$\Delta x_i$ &
Diferencia entre el logaritmo del PBI per cápita relativo de 2019 y el de
2007. \\
\bottomrule
\end{tabular}
\begin{tablenotes}
\footnotesize
\item \textit{Nota.} Las variables de cambio representan la variación
acumulada durante 2007--2019.
\end{tablenotes}
\end{threeparttable}
\end{table}

\subsection{Análisis descriptivo y distributivo}

La primera etapa comprende medidas de tendencia central, dispersión y forma de
la distribución. Se calcularon la media, mediana, valores extremos, desviación
estándar, coeficiente de variación, asimetría, curtosis y el estadístico de
Jarque--Bera.

El coeficiente de variación se utilizó para evaluar la dispersión relativa del
IDH debido a que su media es positiva. En el caso del logaritmo del PBI per
cápita relativo, la media puede ser negativa o cercana a cero; por ello, el
coeficiente de variación no ofrece una interpretación económica adecuada. Para
esta variable se priorizaron la desviación estándar, el rango y la forma de la
distribución.

La evolución distributiva fue complementada mediante estimaciones de densidad
kernel. Para una variable $y$, la densidad estimada se expresa como:

\begin{equation}
\widehat{f}(y)
=
\frac{1}{nh}
\sum_{i=1}^{n}
K\left(
\frac{y-y_i}{h}
\right),
\label{eq:densidad_kernel}
\end{equation}

donde $K(\cdot)$ es la función kernel, $h$ es el ancho de banda y $n$ es el
número de departamentos. La comparación entre las densidades de 2007 y 2019
permite identificar desplazamientos, reducción de dispersión o formación de
agrupamientos.

\subsection{Matrices de pesos espaciales}

La proximidad territorial fue representada mediante tres matrices de pesos. La
especificación principal corresponde a una banda de distancia con ponderación
inversa. Antes de calcular las distancias, los centroides departamentales se
transformaron a un sistema de coordenadas métricas.

Los elementos de la matriz principal se definieron como:

\begin{equation}
w_{ij}
=
\begin{cases}
\dfrac{1}{d_{ij}},
& 0<d_{ij}\leq d_0, \\[0.3cm]
0,
& d_{ij}>d_0,
\end{cases}
\label{eq:matriz_distancia}
\end{equation}

donde $d_{ij}$ representa la distancia entre los centroides de los
departamentos $i$ y $j$, mientras que $d_0$ es el umbral mínimo que garantiza
que todas las unidades territoriales tengan al menos una conexión.

La matriz fue estandarizada por filas:

\begin{equation}
w^{*}_{ij}
=
\frac{w_{ij}}
{\sum_{j=1}^{n}w_{ij}}.
\label{eq:estandarizacion_filas}
\end{equation}

Como pruebas de robustez se construyeron una matriz de contigüidad tipo
\textit{queen}, que considera vecinos a los departamentos que comparten algún
punto de frontera, y una matriz de los cuatro vecinos más próximos
($k=4$). La utilización de matrices alternativas permite evaluar si los
resultados dependen excesivamente de una definición particular de cercanía.

\subsection{Autocorrelación espacial}

La autocorrelación espacial global se evaluó mediante el índice de Moran:

\begin{equation}
I
=
\frac{n}{S_0}
\frac{
\sum_{i=1}^{n}
\sum_{j=1}^{n}
w_{ij}(y_i-\bar{y})(y_j-\bar{y})
}{
\sum_{i=1}^{n}(y_i-\bar{y})^2
},
\label{eq:moran}
\end{equation}

donde:

\begin{equation}
S_0
=
\sum_{i=1}^{n}
\sum_{j=1}^{n}w_{ij}.
\end{equation}

Un valor positivo y estadísticamente significativo indica que los territorios
con valores semejantes tienden a localizarse próximos entre sí. Un valor
negativo significativo representa asociación entre valores territorialmente
diferentes. Cuando el estadístico no es significativo, no se rechaza la
hipótesis de ausencia de autocorrelación espacial global.

El índice se calculó para los niveles iniciales y para los cambios acumulados
del IDH y del PBI per cápita relativo. Los diagramas de dispersión de Moran se
utilizaron como complemento gráfico para identificar la posición relativa de
los departamentos en los cuadrantes alto--alto, bajo--bajo, alto--bajo y
bajo--alto.

\subsection{Convergencia beta global}

La hipótesis de convergencia beta fue estimada inicialmente mediante mínimos
cuadrados ordinarios:

\begin{equation}
\Delta y_i
=
\alpha
+
\beta y_{i,2007}
+
\varepsilon_i,
\label{eq:ols_convergencia}
\end{equation}

donde $y_{i,2007}$ representa el nivel inicial del indicador y $\Delta y_i$ su
cambio acumulado entre 2007 y 2019. Se estimó una ecuación para el IDH y otra
para el logaritmo del PBI real per cápita relativo.

Un coeficiente $\beta<0$ se interpreta como evidencia de convergencia beta,
debido a que las regiones inicialmente rezagadas presentan una variación
posterior mayor. Por el contrario, un coeficiente $\beta>0$ indica divergencia
o persistencia acumulativa de las diferencias iniciales.

Cuando el coeficiente de convergencia es negativo, la velocidad anual se
obtiene mediante:

\begin{equation}
v
=
-\frac{\ln(1+\beta)}{T},
\label{eq:velocidad_convergencia}
\end{equation}

donde $T=12$ representa la duración del periodo 2007--2019.

\subsection{Diagnóstico y modelos espaciales}

Luego de estimar los modelos OLS se aplicaron pruebas de multiplicadores de
Lagrange y sus versiones robustas. Estas pruebas permiten identificar si la
dependencia espacial se concentra en el término de error o en un rezago
espacial.

El modelo de error espacial se representa como:

\begin{align}
\Delta y
&=
X\beta+u, \\
u
&=
\lambda Wu+\varepsilon,
\label{eq:sem}
\end{align}

donde $W$ es la matriz de pesos espaciales, $\lambda$ mide la dependencia entre
los errores regionales y $\varepsilon$ es un término aleatorio independiente.

También se estimó una especificación Durbin de error:

\begin{align}
\Delta y
&=
\alpha
+
\beta y_{2007}
+
\theta Wy_{2007}
+
u, \\
u
&=
\lambda Wu+\varepsilon,
\label{eq:durbin_sem}
\end{align}

donde $Wy_{2007}$ representa el rezago espacial del nivel inicial. Esta
especificación permite evaluar si las condiciones iniciales de los territorios
vecinos están asociadas con el cambio observado en cada departamento.

La selección entre modelos consideró la significancia estadística de los
parámetros espaciales, los resultados de las pruebas LM/RS, el criterio de
información de Akaike y la estabilidad del coeficiente de convergencia. Las
matrices de distancia inversa y contigüidad queen fueron comparadas como parte
del análisis de robustez.

\subsection{Regresión geográficamente ponderada}

La posible variación territorial del proceso de convergencia se examinó
mediante regresiones geográficamente ponderadas. El modelo local se expresa
como:

\begin{equation}
\Delta y_i
=
\alpha(u_i,v_i)
+
\beta(u_i,v_i)y_{i,2007}
+
\varepsilon_i,
\label{eq:gwr}
\end{equation}

donde $(u_i,v_i)$ representa la ubicación geográfica del departamento $i$.
A diferencia del modelo global, la GWR estima un intercepto y un coeficiente
específicos para cada territorio.

Las observaciones cercanas reciben una ponderación mayor que las más distantes.
El ancho de banda fue seleccionado mediante validación cruzada, buscando
minimizar el error de predicción fuera de muestra. Se utilizó un kernel
geográfico continuo para evitar cambios abruptos en la ponderación entre
departamentos.

La interpretación se concentró en el rango de los coeficientes locales, su
distribución cartográfica y la comparación con el coeficiente global. Un rango
local reducido y un ancho de banda próximo a la distancia máxima sugieren que
la relación es esencialmente global. Un rango amplio y un ancho de banda menor
indican una mayor heterogeneidad territorial.

Debido a que la muestra contiene únicamente 24 departamentos, la GWR se
interpreta como una herramienta exploratoria. La amplitud limitada de la
muestra reduce la potencia para detectar variaciones locales y puede producir
estimaciones sensibles al ancho de banda.

\subsection{Regresiones no paramétricas y reproducibilidad}

Finalmente, la relación entre el nivel inicial y el cambio acumulado fue
representada mediante una recta OLS, un suavizamiento LOESS y una regresión
kernel. Estas estimaciones permiten examinar posibles relaciones no lineales
que podrían quedar ocultas en un coeficiente promedio.

El análisis no paramétrico se emplea con fines descriptivos y no se interpreta
como evidencia causal. Su función es comprobar si la dirección de la relación
global se mantiene a lo largo de la distribución o si aparecen cambios de
pendiente en determinados niveles iniciales.

El procesamiento de la información, la construcción de las matrices, las
estimaciones econométricas y la generación de figuras se realizaron en un
cuaderno reproducible. Los datos procesados, el código, las tablas y las
figuras se encuentran organizados en el repositorio digital del proyecto.
